\documentclass[11pt]{article}
\usepackage[margin=1in]{geometry}
\usepackage{amsmath,amssymb,amsthm,mathtools}
\usepackage{graphicx}
\usepackage{float}
\usepackage{enumitem}
\usepackage[colorlinks=true,linkcolor=blue,citecolor=blue,urlcolor=blue]{hyperref}

\newtheorem{theorem}{Theorem}
\newtheorem{lemma}{Lemma}
\newtheorem{proposition}{Proposition}
\theoremstyle{remark}
\newtheorem{remark}{Remark}

\newcommand{\N}{\mathbb N}
\newcommand{\C}{\mathbb C}
\newcommand{\scspec}{\sigma_{\mathrm{sc}}}

\title{Diverging Local Bandwidth Suffices for the\\
Singular-Continuous Three-Limit Lower Bound}
\author{Candidate strong partial result for arXiv:1908.06721}
\date{21 July 2026}

\begin{document}
\maketitle

\begin{center}
\fbox{\parbox{0.91\textwidth}{
\textbf{Status.} Candidate strong partial result, likely valid, awaiting human review.
The theorem removes the source's square-root growth rate, but it does not settle
the bounded-bandwidth case.
}}
\end{center}

\section{Source question}

The source is Matthew J. Colbrook, \emph{Computing Spectral Measures and
Spectral Types} \cite{Colbrook}.  Its identifier is arXiv:1908.06721v3.  For an integer-valued
function $f$ with $f(n)>n$, let $P_n$ be projection onto the first $n$
canonical basis vectors and let
\[
 \Omega_{f,0}=\left\{T=T^*: (I-P_{f(n)})TP_n=0\text{ for every }n\right\}.
\]
Writing $\Xi_{\mathrm{sc}}^{\C}(T)=\scspec(T)$, the source proves in Theorem
5.1 that
\[
 \{\Xi_{\mathrm{sc}}^{\C},\Omega_{f,0},\Lambda_i\}\notin\Delta_3^G,
 \qquad i=1,2,
\]
provided $f(n)-n\geq\sqrt{2n}+1/2$.  In the summary on page 6 it asks whether
this growth condition can be dropped.

\begin{figure}[H]
\centering
\includegraphics[width=\textwidth]{figures/open_problem_crop.png}
\caption{The exact question in the source paper, page 6.}
\end{figure}

The present packet proves the same exclusion whenever $f(n)-n\to\infty$.
It also identifies why the remaining bounded-bandwidth case cannot be reached
by merely improving the enumeration in the source's direct-sum proof.

\section{Proof intuition}

Colbrook reduces a height-three decision problem to singular-continuous
spectrum by taking a direct sum of countably many half-line Schr\"odinger
operators.  The displayed square-root rate enters only when that disjoint
union of rays is placed into one copy of $\ell^2(\N)$ using a triangular
enumeration.

If $h(n)=f(n)-n$ tends to infinity, one can do better.  Wait until the entire
tail has $h\geq k$, and then enumerate the first $k$ rays round-robin.  The gap
between consecutive vertices of any active ray is exactly $k$, so every edge
crossing a cut at $n$ ends no later than $n+h(n)=f(n)$.  Increasing $k$ slowly
gives a bijection with all countably many rays and no prescribed rate of
divergence.

This mechanism is optimal for the direct-sum route.  Across the cut after $n$,
every ray that has already started but is infinite needs a distinct crossing
endpoint among the $h(n)$ sites $n+1,\ldots,f(n)$.  As the number of started
rays tends to infinity, necessarily $h(n)\to\infty$.  Thus bounded bandwidth
requires a genuinely different spectral encoding.

\section{The adaptive enumeration}

We first isolate the combinatorial point.  A \emph{ray enumeration} is a
bijection between $\N$ and $\N\times\N$; for fixed $j$, the pairs $(j,m)$,
$m=1,2,\ldots$, are the successive vertices of ray $j$.

\begin{lemma}[Adaptive round-robin enumeration]\label{lem:enumeration}
Let $f:\N\to\N$ satisfy $f(n)>n$, and put $h(n)=f(n)-n$.  If
$h(n)\to\infty$, there is a ray enumeration with the following property: if
$p<q$ are the positions of two successive vertices of one ray, then
\begin{equation}\label{eq:cut}
 q\leq f(n)\qquad(p\leq n<q).
\end{equation}
\end{lemma}

\begin{proof}
Choose integers
\[
 1=c_1<c_2<c_3<\cdots
\]
recursively so that
\begin{enumerate}[label=(\roman*)]
\item $h(n)\geq k$ for all $n\geq c_k$;
\item $c_{k+1}-c_k$ is a positive multiple of $k$.
\end{enumerate}
This is possible because $h(n)\to\infty$: after choosing $c_k$, move far
enough into the tail where $h\geq k+1$ and then move finitely farther to the
required congruence class modulo $k$.

On the phase $c_k\leq n<c_{k+1}$, label the positions periodically
\[
 1,2,\ldots,k,1,2,\ldots,k,\ldots .
\]
Condition (ii) makes the phase an integral number of cycles.  At phase $k$ the
labels $1,\ldots,k$ occur, and every fixed label occurs in every later phase;
therefore, numbering the successive occurrences of label $j$ by
$m=1,2,\ldots$ gives a bijection $\N\to\N\times\N$.

Within phase $k$, consecutive occurrences of a fixed label are $k$ positions
apart.  At the transition from phase $k$ to phase $k+1$, the last occurrence
of an old label $j\leq k$ is at
\[
 c_{k+1}-k+j-1,
\]
and its first occurrence in the next phase is at $c_{k+1}+j-1$; that gap is
again $k$.  New label $k+1$ has no predecessor.  Consequently, if an edge
$p<q$ starts in phase $k$, then $q-p\leq k$.  For $p\leq n<q$, condition (i)
gives $h(n)\geq k$, and hence
\[
 q-n\leq q-p\leq k\leq h(n).
\]
This is exactly \eqref{eq:cut}.
\end{proof}

\begin{proposition}[Sharpness for disjoint rays]\label{prop:sharp}
Suppose a ray enumeration satisfies \eqref{eq:cut}.  Then
$f(n)-n\to\infty$.
\end{proposition}

\begin{proof}
For a fixed cut $n$, let $A_n$ be the set of rays having at least one vertex
among positions $1,\ldots,n$.  For each $j\in A_n$, take the last vertex of
ray $j$ at or before $n$ and the next vertex of that ray.  The latter is a
crossing endpoint $q_j$ satisfying
\[
 n<q_j\leq f(n).
\]
The endpoints $q_j$ are distinct because the enumeration is a bijection.
Therefore
\[
 |A_n|\leq f(n)-n.
\]
Every one of the countably many rays eventually starts, so for every $K$ one
has $|A_n|\geq K$ for all sufficiently large $n$.  Thus $f(n)-n\to\infty$.
\end{proof}

\section{Strengthened lower bound}

We use the source notation $\Lambda_1$ for matrix-entry evaluations and
$\Lambda_2$ for the augmented evaluation set that also contains the relevant
quadratic matrix data.  The SCI classes and the metric on closed subsets of
$\C$ are exactly those of \cite{Colbrook}.

\begin{theorem}[Diverging-bandwidth lower bound]\label{thm:main}
Let $f:\N\to\N$ satisfy $f(n)>n$ and
\[
 f(n)-n\longrightarrow\infty.
\]
Then, for $i=1,2$,
\[
 \{\Xi_{\mathrm{sc}}^{\C},\Omega_{f,0},\Lambda_i\}
 \notin\Delta_3^G.
\]
Equivalently, no height-two tower of general algorithms computes the
singular-continuous spectrum on $\Omega_{f,0}$.
\end{theorem}

\begin{proof}
We repeat the reduction in the proof of Theorem 5.1 of \cite{Colbrook}, changing
only its enumeration of the direct sum.

Let $a=(a_{k,j})_{k,j\in\N}\in\{0,1\}^{\N\times\N}$ and consider the decision
problem
\[
 \Xi'(a)=1\quad\Longleftrightarrow\quad
 \text{some column of $a$ contains infinitely many ones}.
\]
This problem has general SCI equal to $3$; this is the decision problem used in
the source proof (with the lower bound cited there from
Ben-Artzi--Hansen--Nevanlinna--Seidel \cite{BHNS}).

For every column $j$, define on $\ell^2(\N)$
\[
 H_j=H_0+v_j,\qquad
 v_j(m)=\sum_{k=1}^{\infty}a_{k,j}\,\delta_{m,k!},
\]
where $H_0$ is the free half-line discrete Schr\"odinger operator.  The sparse
potential theorem of Krutikov and Remling \cite{KR}, in precisely the form
quoted as Theorem 3.4 in \cite{Colbrook}, gives
\begin{align*}
 [0,4]&\subset\scspec(H_j) &&\text{if column $j$ has infinitely many ones},\\
 \scspec(H_j)\cap(0,4)&=\varnothing
 &&\text{if column $j$ has only finitely many ones}.
\end{align*}

Set $\widetilde H=\bigoplus_{j=1}^{\infty}H_j$.  Use Lemma
\ref{lem:enumeration} to identify its canonical basis with that of
$\ell^2(\N)$.  The only off-diagonal entries of every $H_j$ join successive
vertices of its ray.  If such an edge has positions $p<q$, then for every cut
$p\leq n<q$, Lemma \ref{lem:enumeration} gives $q\leq f(n)$.  Hence no column
among the first $n$ has a nonzero entry below row $f(n)$, and therefore
\[
 (I-P_{f(n)})\widetilde H P_n=0
 \quad\text{for every }n.
\]
Thus $\widetilde H\in\Omega_{f,0}$.  Its entries are uniformly bounded, so it
is a bounded self-adjoint operator.

The direct-sum spectral decomposition now yields
\[
 \Xi'(a)=1\Longrightarrow [0,4]\subset\scspec(\widetilde H),
 \qquad
 \Xi'(a)=0\Longrightarrow \scspec(\widetilde H)\cap(0,4)=\varnothing.
\]
Every queried matrix entry of $\widetilde H$ depends on at most one entry of
$a$, after applying the fixed enumeration and its inverse.  The augmented
$\Lambda_2$ data are also finite-information functions of $a$, because every
column of $\widetilde H$ has at most three nonzero entries.  Thus the reduction
works for both $\Lambda_1$ and $\Lambda_2$.

If a height-two general tower computed $\scspec(\widetilde H)$, applying the
continuous postprocessing
\[
 F\longmapsto 1-\min\{1,\operatorname{dist}(3,F)\}
\]
and then the discrete-output conversion in the source proof of Theorem 5.1
would give a height-two general tower for $\Xi'$.  This contradicts
$\operatorname{SCI}_G(\Xi')=3$.  The claimed exclusion follows.
\end{proof}

\section{Verification and limitations}

\paragraph{Internal checks.}
The script \texttt{code/check\_round\_robin.py} constructs the phases for the
slowly diverging sample
\[
 h(n)=\lfloor\log_2(n+1)\rfloor+1
\]
and checks every cut inequality in a finite prefix.  This is only a sanity
check; Lemma \ref{lem:enumeration} is the proof.  The most important points for
human verification are (a) the phase-transition gap is $k$, not $k+1$, and
(b) the cut inequality is exactly the matrix condition defining
$\Omega_{f,0}$.

\paragraph{Scope.}
The theorem removes every quantitative growth rate but retains
$f(n)-n\to\infty$.  Proposition \ref{prop:sharp} shows that the latter cannot
be removed from this direct-sum/ray-enumeration proof.  The original question
remains open for bounded or non-diverging local bandwidth, notably
$f(n)=n+1$.

\paragraph{Novelty check.}
The run's lightweight registry, solution, attempt, and proof-gap indexes were
searched for arXiv:1908.06721, the title, singular-continuous spectrum, SCI,
and bandwidth terminology; no duplicate was found.  A bounded web/arXiv search
on 21 July 2026 used the exact source sentence and the phrases
``singular continuous spectrum'', ``Delta 3'', ``SCI'', ``bandwidth'', and
``Colbrook''.  It found the source and later numerical spectral-measure papers
but no separate settlement of this hypothesis.  The novelty confidence is
moderate rather than definitive because this was not an exhaustive citation
review.

\paragraph{Recommendation.}
Promote as a strong partial result if a reviewer confirms Lemma
\ref{lem:enumeration} and the finite-information simulation for $\Lambda_2$.
Do not describe it as a full answer to the source question.

\begin{thebibliography}{9}

\bibitem{Colbrook}
M.~J. Colbrook,
\emph{Computing Spectral Measures and Spectral Types},
Communications in Mathematical Physics 384 (2021), 433--501;
arXiv:1908.06721v3.

\bibitem{KR}
D.~Krutikov and C.~Remling,
\emph{Schr\"odinger operators with sparse potentials: asymptotics of the
Fourier transform of the spectral measure},
Communications in Mathematical Physics 223 (2001), 509--532.

\bibitem{BHNS}
J.~Ben-Artzi, M.~J. Colbrook, A.~C. Hansen, O.~Nevanlinna, and M.~Seidel,
\emph{Computing Spectra---On the Solvability Complexity Index hierarchy and
towers of algorithms},
arXiv:1508.03280v5; see also the exact decision-problem citation in
\cite{Colbrook}.

\end{thebibliography}

\end{document}
